\section*{Hoofdstuk \ref{ch:b2:aromatic-substitution} --- Aromaticiteit en elektrofiele aromatische substitutie}

\begin{solution}{exo:b2:aromatic-substitution:1}
\ce{HNO3 + 2H2SO4 -> NO2+ + H3O+ + 2HSO4-}: salpeterzuur wordt geprotoneerd en verliest dan water.
\end{solution}

\begin{solution}{exo:b2:aromatic-substitution:2}
Tolueen geeft 2-broomtolueen en 4-broomtolueen: de methylgroep stuurt naar de ortho- en
paraposities. Nitrobenzeen geeft 3-broomnitrobenzeen, trager: de nitrogroep
stuurt naar meta.
\end{solution}

\begin{solution}{exo:b2:aromatic-substitution:3}
\ce{-CH3}: activerend, o/p. \ce{-OH}: sterk activerend, o/p. \ce{-Cl}: desactiverend, o/p.
\ce{-NO2}: sterk desactiverend, meta. \ce{-COCH3}: desactiverend, meta. \ce{-NHCOCH3}:
activerend (matig), o/p.
\end{solution}

\begin{solution}{exo:b2:aromatic-substitution:4}
Acetofenon (1-fenylethaan-1-on), \ce{C6H5COCH3}, via het \omterm{def:b2:aromatic-substitution:friedel-crafts}{acyliumion} \ce{CH3CO+}.
\end{solution}

\begin{solution}{exo:b2:aromatic-substitution:5}
Het primaire carbokation (of zijn complex) legt om door een hydrideverschuiving tot het secundaire
isopropylkation, dat de ring alkyleert. Propylbenzeen: \omterm{def:b2:aromatic-substitution:friedel-crafts}{friedel-craftsacylering} met
propanoylchloride (geen omlegging), daarna reductie van de ketongroep \ce{C=O} tot \ce{CH2}.
\end{solution}

\begin{solution}{exo:b2:aromatic-substitution:6}
De groep \ce{-OH} is een sterke mesomere donor: de ring is zo geactiveerd dat moleculair broom
een voldoende sterk elektrofiel is, en alle drie de ortho- en paraposities worden gesubstitueerd.
\end{solution}

\begin{solution}{exo:b2:aromatic-substitution:7}
3-Broomnitrobenzeen: nitreren, dan bromeren (nitro stuurt naar meta). 4-Broomnitrobenzeen:
bromeren, dan nitreren (broom stuurt naar ortho/para), en het para-isomeer scheiden van het
ortho-isomeer.
\end{solution}

\begin{solution}{exo:b2:aromatic-substitution:8}
Het nitreermengsel oxideert het zeer elektronenrijke aniline en protoneert zijn stikstofatoom:
\ce{-NH3+} is een meta-sturende, \omterm{def:b2:aromatic-substitution:directing}{desactiverende groep}. In aceetanilide wordt het vrije elektronenpaar van stikstof
gedeeld met de carbonylgroep: minder basisch en minder activerend, wordt het niet geprotoneerd of geoxideerd,
en het stuurt nog steeds naar ortho/para, vooral naar para om sterische redenen.
\end{solution}

\begin{solution}{exo:b2:aromatic-substitution:9}
Sulfoneer tolueen (vooral para, de omvangrijke groep vermijdt ortho); bromeer: het broom komt
ortho ten opzichte van de methylgroep (de parapositie is geblokkeerd, en beide groepen begunstigen die positie);
verwijder de groep \ce{SO3H} door te verhitten in verdund zuur.
\end{solution}

\begin{solution}{exo:b2:aromatic-substitution:10}
Beide groepen sturen naar ortho/para; hun paraposities worden door elkaar bezet. Methoxy, de
sterkere activator, wint: nitrering ortho ten opzichte van \ce{-OCH3} (positie 2), wat
4-methyl-2-nitroanisool geeft.
\end{solution}

\begin{solution}{exo:b2:aromatic-substitution:11}
Elke nitrogroep desactiveert de ring sterk: de tweede nitrering vraagt strengere omstandigheden
dan de eerste, de derde (op een ring met twee nitrogroepen en alleen de zwak activerende
methylgroep) nog veel strengere.
\end{solution}

\begin{solution}{exo:b2:aromatic-substitution:12}
4-Nitrobenzoëzuur: nitreer tolueen, scheid het para-isomeer af, oxideer de methylgroep tot
\ce{-COOH} (heet permanganaat). 3-Nitrobenzoëzuur: oxideer tolueen eerst tot benzoëzuur, en
nitreer daarna (\ce{-COOH} stuurt naar meta).
\end{solution}

\begin{solution}{pb:b2:aromatic-substitution:1}
\textbf{1.} \ce{HNO3 + 2H2SO4 -> NO2+ + H3O+ + 2HSO4-}.
\textbf{2.} \ce{O=N+=O}: lineair, iso-elektronisch met \ce{CO2}.
\textbf{3.} Het stikstofatoom bindt aan het ringkoolstofatoom para ten opzichte van \ce{CH3}, dat tetraëdrisch wordt
(en H en \ce{NO2} draagt).
\textbf{4.} De positieve lading op de twee koolstofatomen ortho ten opzichte van het aangevallen koolstofatoom en op het
koolstofatoom para ertegenover, dat de methylgroep draagt.
\textbf{5.} De eerste, de vorming van het \omterm{def:b2:aromatic-substitution:seAr}{whelandintermediair}.
\textbf{6.} Een base uit het medium: \ce{HSO4-} of water.
\textbf{7.} Twee ortho-, twee meta-, één parapositie.
\textbf{8.} 2 : 2 : 1, dus 40~\% ortho, 40~\% meta, 20~\% para.
\textbf{9.} Bij meta-aanval staat de positieve lading nooit op het koolstofatoom dat de methylgroep draagt; geen enkele
structuur krijgt hulp van de methylgroep.
\textbf{10.} Eén resonantiestructuur heeft de lading op het koolstofatoom dat \ce{CH3} draagt: een tertiair
carbokation, gestabiliseerd door het inductieve effect en de \omterm{def:b2:fragment-orbitals:hyperconjugation}{hyperconjugatie} van de methylgroep.
\textbf{11.} Ortho en para worden begunstigd (96~\% tegenover 60~\% bij toeval); meta is bijna afwezig.
\textbf{12.} Per positie: ortho $58/2 = 29~\%$, para 38~\%: para is de reactiefste positie,
want de orthoposities worden een beetje gehinderd door de methylgroep.
\textbf{13.} Sneller: de methylgroep activeert de ring.
\textbf{14.} 4-Nitrotolueen (smelt bij \qty{52}{\degreeCelsius}).
\textbf{15.} Koel het ruwe mengsel af: het para-isomeer kristalliseert en wordt afgefiltreerd; de
vloeistof houdt het ortho- en het meta-isomeer.
\textbf{16.} Een symmetrisch molecuul stapelt beter in een kristal: meer interacties per volume,
een hoger smeltpunt.
\textbf{17.} Door \omterm{def:b2:liquid-vapour-diagrams:fractional-distillation}{gefractioneerde destillatie} onder verlaagde druk, of door chromatografie; of door
verder af te koelen (het meta-isomeer smelt bij \qty{16}{\degreeCelsius}).
\textbf{18.} De isomeren hebben dicht bij elkaar liggende kookpunten (zelfde formule, vergelijkbare polariteit).
\textbf{19.} \ce{C7H8}: \qty{92.14}{g/mol}; \ce{C7H7NO2}: \qty{137.14}{g/mol}.
\textbf{20.} $100/92.14 = \qty{1.085}{mol}$ tolueen; $0.90 \times 1.085 = \qty{0.977}{mol}$
nitrotoluenen.
\textbf{21.} Totaal $0.977 \times 137.14 = \qty{134.0}{g}$: ortho \qty{77.7}{g}, meta
\qty{5.4}{g}, para \qty{50.9}{g}.
\textbf{22.} Een tweede nitrering (dinitrotoluenen), en oxidatie van de methylgroep.
\textbf{23.} Oxideer de methylgroep tot \ce{-COOH} met heet basisch permanganaat, en zuur daarna aan.
\textbf{24.} De run geeft $\boldsymbol{\approx \qty{51}{g}}$ 4-nitrotolueen.
\end{solution}
